\section{Aqueous Reactions, Redox Balancing, and Gases}

\subsection{Ctrl+F keywords: molarity, dilution, net ionic, precipitation, redox, oxidation number, gas law, partial pressure, ideal gas, gas stoichiometry, pressure tank, bundfald}
Use for solution concentrations, net ionic equations, precipitate decisions, half-reaction balancing, gas state changes, gas mixtures, and reaction amounts involving gases.

\subsection{Symbols and notation}
\referencetable{tab:aqueous:symbols}{Symbols used for aqueous reactions, redox, and gases}
\begin{tabular}{@{}>{$}l<{$} p{10.8cm}@{}}
\toprule
\text{Symbol} & \text{Meaning and usage}\\
\midrule
c,\ n,\ V & Molar concentration, amount of substance, and solution or gas volume.\\
c_1,V_1,c_2,V_2 & Concentration and volume before and after dilution.\\
Q_{\mathrm{sp}},\ K_{\mathrm{sp}} & Ion product at current concentrations and solubility-product constant.\\
p,\ T,\ R & Gas pressure, absolute temperature, and gas constant.\\
x_i,\ p_i,\ p_{\mathrm{tot}} & Mole fraction, partial pressure of gas \(i\), and total pressure.\\
\rho,\ M & Gas density and molar mass.\\
\ce{e-},\ \ce{H+},\ \ce{OH-} & Species used in half-reaction balancing.\\
\ce{(aq)},\ \ce{(s)} & Dissolved aqueous species and solid precipitate.\\
\bottomrule
\end{tabular}

\subsection{Solutions and aqueous reactions}

\subsubsection{Molarity, dilution, actual ion concentration}
\begin{equation}\label{eq:aqueous:concentration-dilution}
c=\frac{n}{V}, \qquad n=cV, \qquad c_1V_1=c_2V_2
\end{equation}
Use liters for \(V\). For a completely dissociated strong electrolyte, obtain an ion concentration by multiplying \(c\) by its coefficient; for example, \(\ce{CaCl2(aq) -> Ca^2+(aq) + 2Cl-(aq)}\) gives \([\ce{Cl-}]=2c\).

\subsubsection{Complete dissolution and total ion amount}
\begin{equation}\label{eq:aqueous:complete-dissolution-ions}
\mathrm{A}_a\mathrm{B}_b(aq)\longrightarrow a\,\mathrm{A}(aq)+b\,\mathrm{B}(aq), \qquad
n_{\mathrm{ions,total}}=(a+b)n_{\mathrm{salt}}
\end{equation}
For a strong soluble electrolyte, each formula unit produces the number of dissolved ions shown by the balanced dissociation coefficients.

\subsection{Exam recipe: moles of ions from complete dissolution}
\textbf{Use when:} a soluble ionic compound dissolves completely and the task asks for total moles of ions or the amount of one ion.

\textbf{Given / Find:} Given the amount or mass of a soluble salt; find \(n_{\mathrm{ions,total}}\) or \(n\) of a specified ion.

\textbf{Required references:} Equations \eqref{eq:aqueous:complete-dissolution-ions} and \eqref{eq:aqueous:concentration-dilution}; use Equation \eqref{eq:atoms:mole-mass-particles} if a mass is given.

\textbf{Steps:}
\begin{enumerate}
    \item Write the complete dissociation equation and balance the ionic coefficients; for example, one \(\ce{Al2(SO4)3}\) formula unit produces \(2+3=5\) ions.
    \item Convert a stated mass or solution concentration to moles of salt before applying the dissociation coefficients.
    \item Multiply salt moles by the requested ion coefficient, or by the sum of all ion coefficients for total moles of ions, using \eqref{eq:aqueous:complete-dissolution-ions}.
\end{enumerate}

\textbf{Checks:} Total charge of ions produced from one formula unit is zero; total ion moles exceed salt moles when more than one ion is formed.

\textbf{Common traps:} counting distinct ion types instead of stoichiometric ion amounts; forgetting parentheses in salts such as \(\ce{Al2(SO4)3}\).

\subsubsection{Net ionic equation and spectator ions}
\begin{equation}\label{eq:aqueous:net-ionic-example}
\ce{Pb^2+(aq) + 2I-(aq) -> PbI2(s)}
\end{equation}
Remove ions that occur unchanged on both sides. A net ionic equation conserves atoms and total charge.

\subsubsection{Precipitation recognition}
\begin{equation}\label{eq:aqueous:precipitation-test}
Q_{\mathrm{sp}}>K_{\mathrm{sp}}\Rightarrow\text{precipitate}, \qquad
Q_{\mathrm{sp}}<K_{\mathrm{sp}}\Rightarrow\text{no precipitate yet}
\end{equation}
\referencetable{tab:aqueous:precipitation-language}{Recognition cues for precipitation questions}
\begin{tabular}{@{}>{\raggedright\arraybackslash}p{4.2cm}>{\raggedright\arraybackslash}p{8.0cm}@{}}
\toprule
Problem wording & Immediate action\\
\midrule
mixing ionic solutions; bundfald & Form candidate products and test whether one is insoluble.\\
\(K_{\mathrm{sp}}\) or sparingly soluble & Compute post-mixing ion concentrations and compare \(Q_{\mathrm{sp}}\) with \(K_{\mathrm{sp}}\).\\
net ionic equation & Dissociate soluble strong electrolytes and cancel spectator ions.\\
\bottomrule
\end{tabular}

\referencetable{tab:aqueous:solubility-rules}{Common solubility rules for deciding candidate precipitates when no \(K_{\mathrm{sp}}\) data are supplied}
\begin{tabular}{@{}>{\raggedright\arraybackslash}p{4.2cm}>{\raggedright\arraybackslash}p{8.0cm}@{}}
\toprule
Ions in compound & Rapid solubility decision\\
\midrule
\(\ce{NO3-}\), group 1 cations, \(\ce{NH4+}\) & Treat as soluble; these usually remain spectator ions.\\
\(\ce{Cl-}\), \(\ce{Br-}\), \(\ce{I-}\) & Usually soluble; treat compounds with \(\ce{Ag+}\) or \(\ce{Pb^2+}\) as precipitate candidates.\\
\(\ce{SO4^2-}\) & Usually soluble; treat compounds with \(\ce{Ba^2+}\) or \(\ce{Pb^2+}\) as precipitate candidates.\\
\(\ce{CO3^2-}\), \(\ce{PO4^3-}\), \(\ce{OH-}\) & Usually insoluble unless paired with group 1 cations or \(\ce{NH4+}\).\\
\bottomrule
\end{tabular}

\subsection{Exam recipe: net ionic precipitation}
\textbf{Use when:} two aqueous ionic solutions are mixed and the question asks for a precipitate or net ionic equation.

\textbf{Given / Find:} Given dissolved compounds, concentrations, and possibly \(K_{\mathrm{sp}}\); find whether a solid forms and write its net ionic equation.

\textbf{Required references:} Equations \eqref{eq:aqueous:concentration-dilution}, \eqref{eq:aqueous:net-ionic-example}, \eqref{eq:aqueous:precipitation-test}, and Tables \ref{tab:aqueous:precipitation-language} and \ref{tab:aqueous:solubility-rules}.

\textbf{Steps:}
\begin{enumerate}
    \item Dissociate each soluble strong electrolyte into ions and obtain concentrations after mixing with \(c=n/V\) from \eqref{eq:aqueous:concentration-dilution} when numerical data are given.
    \item Pair the available cations and anions to write candidate neutral solid formulas by charge balance.
    \item Decide the precipitate branch: when no numerical solubility data are supplied, classify each candidate using Table \ref{tab:aqueous:solubility-rules}; when \(K_{\mathrm{sp}}\) is supplied, calculate \(Q_{\mathrm{sp}}\) from mixed ion concentrations and apply \eqref{eq:aqueous:precipitation-test}.
    \item Write only the precipitating ions and the solid product, following the pattern in \eqref{eq:aqueous:net-ionic-example}, then choose coefficients that conserve atoms and charge.
\end{enumerate}

\textbf{Checks:} No unchanged spectator ion remains; phases and charges are shown; both atom count and charge balance.

\textbf{Common traps:} Using concentrations before dilution by mixing; leaving nitrate or alkali metal ions in the net equation; writing an uncharged ionic product with incorrect subscripts.

\subsubsection{Gravimetric concentration from a precipitate}
\begin{equation}\label{eq:aqueous:gravimetric-precipitation}
n_{\mathrm{precipitate}}=\frac{m_{\mathrm{precipitate}}}{M_{\mathrm{precipitate}}},
\qquad
n_{\mathrm{analyte}}=n_{\mathrm{precipitate}}
\frac{\nu_{\mathrm{analyte}}}{\nu_{\mathrm{precipitate}}},
\qquad
c_{\mathrm{analyte}}=\frac{n_{\mathrm{analyte}}}{V_{\mathrm{sample}}}
\end{equation}
The coefficient ratio is taken from the balanced net ionic precipitation reaction.

\subsection{Exam recipe: concentration from precipitate mass}
\textbf{Use when:} an ion is precipitated quantitatively, the mass of solid is measured, and the original solution concentration is requested.

\textbf{Given / Find:} Given a precipitate mass, its identity, the sampled solution volume, and the analyte ion; find the analyte concentration.

\textbf{Required references:} Equation~\eqref{eq:aqueous:gravimetric-precipitation}, Equation~\eqref{eq:atoms:mole-mass-particles}, and the net ionic reaction method in Equation~\eqref{eq:aqueous:net-ionic-example}.

\textbf{Steps:}
\begin{enumerate}
    \item Write the balanced net ionic reaction linking the analyte ion to the measured solid; its coefficients define the mole ratio.
    \item Calculate the precipitate molar mass, then convert measured solid mass to moles using the first expression in Equation~\eqref{eq:aqueous:gravimetric-precipitation}.
    \item Convert precipitate moles to analyte moles using the coefficient ratio in the second expression of Equation~\eqref{eq:aqueous:gravimetric-precipitation}.
    \item Divide analyte moles by the original sampled solution volume in liters using the final expression of Equation~\eqref{eq:aqueous:gravimetric-precipitation}.
\end{enumerate}

\textbf{Checks:} For \(\ce{Ag+ + Cl- -> AgCl(s)}\), precipitate and chloride moles are equal; concentration units are \(\mathrm{mol\,L^{-1}}\).

\textbf{Common traps:} dividing by total reagent-plus-sample volume rather than the original sample volume; using mass ratios in place of the balanced mole ratio.

\subsection{Redox balancing}

\subsubsection{Oxidation, reduction, and electron cancellation}
Oxidation produces electrons and increases oxidation state; reduction consumes electrons and decreases oxidation state. The final equation contains no free electrons.

\subsubsection{Half-reaction sequence in acid and base}
\begin{equation}\label{eq:aqueous:redox-half-reaction-sequence}
\begin{aligned}
&\text{balance atoms except H and O};\quad
\text{balance O with }\ce{H2O};\\
&\text{balance H with }\ce{H+};\quad
\text{balance charge with }\ce{e-};\\
&\text{in base, add }\ce{OH-}\text{ to both sides for every }\ce{H+}\text{ and cancel }\ce{H2O}.
\end{aligned}
\end{equation}

\subsubsection{Oxidation-state formula:}
\begin{equation}\label{eq:aqueous:oxidation-state-sum}
\sum_i n_i\,\mathrm{OS}_i = q_{\text{species}}
\end{equation}
where \(n_i\) is the number of atoms of element \(i\), \(\mathrm{OS}_i\) is its oxidation state, and \(q_{\text{species}}\) is the total charge of the molecule or ion. Assign known values first, usually \(\ce{O}=-II\), \(\ce{H}=+I\), free elements \(=0\), and monatomic ions \(=\) their ion charge; set the unknown oxidation state to \(x\) and solve.

\subsection{Exam recipe: redox balancing}
\textbf{Use when:} an oxidation--reduction equation must be balanced in acidic or basic aqueous solution.

\textbf{Given / Find:} Given reactant and product species plus acidic or basic medium; find the balanced overall equation.

\textbf{Required references:} Oxidation-state rules in Table~\ref{tab:electrochem:oxidation-rules}; half-reaction procedure \eqref{eq:aqueous:redox-half-reaction-sequence}.

\textbf{Steps:}
\begin{enumerate}
    \item Identify oxidation states with \eqref{eq:aqueous:oxidation-state-sum} for the atoms that change: multiply each atom count by its oxidation state, set the sum equal to the species charge, and solve any unknown \(x\).
    \item Separate oxidation and reduction half-reactions: oxidation state increases for oxidation and decreases for reduction.
    \item For each half-reaction, balance non-\(\ce{H}\), non-\(\ce{O}\) atoms, then add \(\ce{H2O}\), \(\ce{H+}\), and \(\ce{e-}\) in the order prescribed by \eqref{eq:aqueous:redox-half-reaction-sequence}.
    \item Multiply each balanced half-reaction by the smallest integers that give equal electron counts, then add and cancel identical species.
    \item For basic solution, apply the final conversion step in \eqref{eq:aqueous:redox-half-reaction-sequence} and cancel excess water.
\end{enumerate}

\textbf{Checks:} Total atoms and total charge agree on both sides; electrons do not occur in the final equation; no \(\ce{H+}\) remains in a final basic-medium equation.

\textbf{Common traps:} Balancing charge before oxygen and hydrogen; converting to base before the acidic half-reaction is balanced; failing to cancel water.

\subsection{Gases}

\subsubsection{Ideal gas law and fixed-amount state change}
\begin{equation}\label{eq:aqueous:ideal-gas-law}
pV=nRT
\end{equation}
\begin{equation}\label{eq:aqueous:combined-gas-law}
\frac{p_1V_1}{T_1}=\frac{p_2V_2}{T_2}\qquad(n\ \text{constant})
\end{equation}
Temperatures must be in kelvin. Choose \(R\) to match the pressure and volume units.

\subsubsection{Partial pressure and mole fraction}
\begin{equation}\label{eq:aqueous:dalton-law}
p_i=x_ip_{\mathrm{tot}}, \qquad x_i=\frac{n_i}{n_{\mathrm{tot}}}, \qquad p_{\mathrm{tot}}=\sum_i p_i
\end{equation}
For an ideal gas mixture, mole fraction equals pressure fraction.

\subsubsection{Gas density and molar mass}
\begin{equation}\label{eq:aqueous:gas-density}
\rho=\frac{m}{V}=\frac{pM}{RT}, \qquad M=\frac{\rho RT}{p}
\end{equation}
Use \eqref{eq:aqueous:gas-density} when density, pressure, and temperature identify an unknown gas.

\subsection{Exam recipe: molar mass and molecular formula from gas data}
\textbf{Use when:} an unknown gaseous compound is specified by mass and gas-state data and its molar mass or molecular formula is requested.

\textbf{Given / Find:} Given gas mass, pressure, volume, temperature, and optionally an empirical formula; find \(M\) and/or the molecular formula.

\textbf{Required references:} Equations~\eqref{eq:aqueous:ideal-gas-law} and \eqref{eq:aqueous:gas-density}; use Equation~\eqref{eq:atoms:molecular-factor} when an empirical formula is given.

\textbf{Steps:}
\begin{enumerate}
    \item Convert temperature to kelvin and use consistent pressure-volume units with the chosen gas constant.
    \item Calculate gas moles as \(n=pV/(RT)\) from Equation~\eqref{eq:aqueous:ideal-gas-law}, then calculate \(M=m/n\); equivalently calculate \(\rho=m/V\) and use Equation~\eqref{eq:aqueous:gas-density}.
    \item If an empirical formula is supplied or obtained, calculate its formula mass and determine the integer factor \(k=M/M_{\mathrm{empirical}}\) by Equation~\eqref{eq:atoms:molecular-factor}.
    \item Multiply every empirical-formula subscript by the nearest justified integer \(k\) to state the molecular formula.
\end{enumerate}

\textbf{Checks:} \(M\) must be positive and close to an integer multiple of the empirical-formula mass when a molecular formula is requested.

\textbf{Common traps:} using Celsius; treating gas volume as molar mass information without pressure and temperature; rounding \(k\) before checking measurement precision.

\subsubsection{Effusion through a small opening, Graham's law}
\begin{equation}\label{eq:aqueous:graham-effusion}
\frac{r_1}{r_2}=\sqrt{\frac{M_2}{M_1}}
\end{equation}
At the same temperature, the lighter gas effuses faster because its molar mass is smaller.

\subsection{Exam recipe: gas law state change}
\textbf{Use when:} a fixed amount of gas changes pressure, volume, or temperature without a chemical reaction, such as a sealed tank or piston problem.

\textbf{Given / Find:} Given two gas states with one unknown variable; find final pressure, volume, or temperature.

\textbf{Required references:} Combined gas law \eqref{eq:aqueous:combined-gas-law}; use \eqref{eq:aqueous:ideal-gas-law} only when moles must first be determined.

\textbf{Steps:}
\begin{enumerate}
    \item Confirm that the amount of gas is constant; otherwise calculate or track \(n\) with \eqref{eq:aqueous:ideal-gas-law} instead of applying \eqref{eq:aqueous:combined-gas-law}.
    \item Convert each temperature to kelvin and put pressure and volume in matching units on both sides.
    \item Substitute the two states into \eqref{eq:aqueous:combined-gas-law} and algebraically isolate the requested variable.
\end{enumerate}

\textbf{Checks:} At constant volume, increasing absolute temperature increases pressure; at constant temperature, increasing volume decreases pressure.

\textbf{Common traps:} Substituting Celsius; applying the combined gas law when gas is added, removed, or consumed.

\subsection{Exam recipe: gas stoichiometry}
\textbf{Use when:} a reacting gas is specified by pressure, volume, temperature, or partial pressure, including a closed heated tank in which gases react before a final pressure is requested.

\textbf{Given / Find:} Given a balanced reaction and gas-state data; find reacting moles, product mass, product gas volume, or final gas pressure.

\textbf{Required references:} Equations \eqref{eq:aqueous:ideal-gas-law}, \eqref{eq:aqueous:dalton-law}, and \eqref{eq:atoms:stoichiometric-conversion}.

\textbf{Steps:}
\begin{enumerate}
    \item If the gas is in a mixture, calculate its partial pressure with \eqref{eq:aqueous:dalton-law}; otherwise use the stated gas pressure.
    \item Convert the reacting gas data to moles by rearranging \eqref{eq:aqueous:ideal-gas-law} as \(n=pV/(RT)\), using kelvin and compatible units.
    \item Write and balance the reaction, then apply the coefficient ratio exactly as in \eqref{eq:atoms:stoichiometric-conversion}.
    \item For a final reacting-gas mixture, subtract consumed reactants, add gaseous products, sum remaining gas moles, and use \eqref{eq:aqueous:ideal-gas-law} at the final \(T\) and \(V\); otherwise convert product moles directly to the requested mass, volume, or pressure.
\end{enumerate}

\textbf{Checks:} The pressure used is the reacting gas pressure; balanced coefficients relate moles, not liters unless state conditions are identical.

\textbf{Common traps:} Using total pressure for one component; applying gas volume ratios at different temperatures or pressures; using Celsius in \(pV=nRT\).

\subsection{Exam recipe: gas effusion from a punctured container}
\textbf{Use when:} a container is punctured, gases escape through a small hole, or the task compares which gas remains or escapes faster.

\textbf{Given / Find:} Given two gas identities or molar masses in one container; find relative effusion rate or enrichment of the remaining gas.

\textbf{Required references:} Equation \eqref{eq:aqueous:graham-effusion}.

\textbf{Steps:}
\begin{enumerate}
    \item Determine each molar mass from its formula.
    \item Compare rates using \eqref{eq:aqueous:graham-effusion}: the gas with smaller \(M\) has larger effusion rate.
    \item For a punctured mixture, state that the faster-effusing gas becomes depleted more quickly in the container, so the remaining gas mixture becomes richer in the heavier gas.
\end{enumerate}

\textbf{Checks:} If \(M_1<M_2\), then \(r_1/r_2>1\); identical molar masses give equal rates.

\textbf{Common traps:} using molecular size intuition instead of molar mass; concluding that the faster gas becomes enriched inside the leaking container.
